\section{Discussion}
\label{sec:discussion}

\paragraph{GLOF contribution to the AMS.}
Under daily maximum discharge, the outburst produced the annual maximum in 4 of 8
GLOF years (2012, 2013, 2016, 2018). In the remaining four (2011, 2014, 2015, 2017),
independent storm floods recorded higher instantaneous peaks, including a
$\hat{T} \approx 127$-year rainfall event in 2014 ($Q = 40.48$\,\si{m^3\per s}),
two approximately 17--20-year events (2011, 2017), and one minor storm year
(2015, $\hat{T} \approx 2$\,yr). Under daily mean discharge, GLOFs dominated in
6 of 8 years, because their sustained multi-hour outflows elevate the 24-hour average
above brief convective storm peaks. The choice of aggregation statistic therefore
fundamentally shapes the apparent GLOF influence on the AMS: daily maxima are the more
relevant basis for short-duration hydraulic design, while daily means better characterise
flood volume relevant to inundation duration or reservoir routing.

\paragraph{Statistical signal driven by the 2018 event.}
The 2018 outburst ($Q = 56.90$\,\si{m^3\per s}, $\hat{T} \approx 789$\,yr under the
reference GEV) is the primary reason the GLOF-inclusive values fail the KS test. A nearly 800-year event
in an 8-year window is statistically improbable but not impossible ($p \approx 1\,\%$).
The causal attribution to the GLOF rests on the masking experiment: after the $\pm 5$-day
window, 2018's annual maximum collapses from \SI{56.90}{m^3\per s} to
\SI{14.28}{m^3\per s} ($\hat{T} \approx 1.5$\,yr), confirming that the extreme peak was
attributable to the outburst, not concurrent meteorological forcing. With only $n = 8$
test observations and a critical value $C_{0.05} = 0.455$ that is itself sensitive to a
single outlier, the formal KS test carries limited inferential weight in isolation; the
masking experiment provides more compelling mechanistic evidence.

\paragraph{Hypothesis 1 confirmed.}
The KS test rejects $H_0^{(1)}$ in both datasets (BAFU~2219: $p = 0.002$; GRDC: $p < 0.001$).
Rank-Sum comparison~5a, which directly compares the pre-GLOF reference against the
GLOF-inclusive values, corroborates this (BAFU~2219: $p = 0.002$; GRDC: $p < 0.001$).
Comparison~5c, using all non-GLOF years as the reference group, confirms the rejection
with consistent $p$-values. The stronger rejection for GRDC is consistent with GLOFs
dominating 6 of 8 years under daily means versus 4 of 8 under daily maxima. The
full AMS GEV overstates the 100-year design flood by \SI{15.0}{\percent} for BAFU~2219
and \SI{15.3}{\percent} for GRDC relative to the GLOF-masked estimate, driven by an
upward shift in shape parameter ($\hat{\xi}: 0.160 \to 0.251$ for BAFU~2219). The
inflation is dominated by the 2018 event rather than being a systematic feature of all
eight GLOFs.

\paragraph{Hypothesis 2 supported.}
The KS test fails to reject $H_0^{(2)}$ in both datasets (BAFU~2219: $p = 0.195$;
GRDC: $p = 0.321$). Rank-Sum comparison~5b, which directly tests the same null
(pre-GLOF reference vs.\ GLOF-masked), also fails to reject (BAFU~2219: $p = 0.172$;
GRDC: $p = 0.233$). Comparison~5d, using all non-GLOF years as the reference group,
corroborates (BAFU~2219: $p = 0.182$; GRDC: $p = 0.226$). Together these confirm that
once GLOF contributions are removed, the natural-flood regime during 2011--2018 is
statistically indistinguishable from the pre-GLOF reference. The convergence of results
across two different discharge statistics strengthens this conclusion.

Nonetheless, Figs.~\ref{fig:gev_three} and~\ref{fig:pp} reveal a consistent upward
tendency in the GLOF-masked peaks: the masked GEV curve lies above the reference across
all return periods (Fig.~\ref{fig:gev_three}), and all eight GLOF-masked annual maxima
plot above the reference diagonal in the probability plot (Fig.~\ref{fig:pp}). This
pattern could reflect gradual hydroclimatic change during 2011--2018 --- for instance,
increasing precipitation intensity under a warmer climate --- superimposed on the natural
reference variability. However, with only $n = 8$ observations the KS test ($p = 0.195$)
and Rank-Sum comparison~5b ($p = 0.172$) cannot distinguish this systematic elevation
from sampling variability, and $H_0^{(2)}$ is not rejected at $\alpha = 0.05$. A longer
post-bypass record is needed to determine whether this upward tendency represents a
persistent hydroclimatic trend.

\paragraph{Implications of the bypass intervention.}
The bypass channel constructed in 2019 \citep{GeoTest2019} has ended the uncontrolled
GLOF episode. The GLOF-masked or pre-GLOF reference GEV is now the appropriate design
distribution for natural flood hazard at Lenk. Naive inclusion of unmasked peaks
overstates the 100-year design flood by 15\,\% and the 200-year design flood by
18--20\,\%, which would lead to unnecessarily conservative standards in any post-bypass
hazard reassessment.

\paragraph{Limitations.}
The principal limitation is sample size: only $n = 8$ GLOF-period years constrain the
test groups, making both tests sensitive to a single outlier. Failure to reject Hypothesis~2 cannot
be interpreted as confirmation of stationarity. Future work should extend the analysis as
post-bypass data accumulate, apply mixture models that explicitly separate the GLOF and
natural-flood populations \citep{Kite1977}, and monitor how lake drainage dynamics evolve
as Plaine Morte continues to lose mass.
